% Standalone appendix document - compiles on its own:
%   latexmk -pdf 01-network-architecture.tex
\documentclass[a4paper,11pt]{article}
\usepackage[english]{babel}
\usepackage{../../appendix-style}

\title{3.1 Network Architecture}
\author{SW4PRJ4 -- Group 3}
\date{\today}

\begin{document}
\maketitle

\begin{versionhistory}
  1.0 & 16-09-2026 & Group 3 & Local and production setup with docker compose \\
\end{versionhistory}

\section{Overview}
PrepMeUp runs as a set of Docker containers started with docker compose. There are two setups: one for local development (\texttt{docker-compose.yml}) and one for production on our server with Coolify (\texttt{docker-compose.prod.yml}). In both, all containers are on the same Docker network, \texttt{prepmeup-net}, and find each other by service name, for example \texttt{db} and \texttt{keycloak}. Why we chose Keycloak and Coolify is described in appendix 2.1.

The two clients are not containers in compose. The customer app runs on the phone, and the admin panel runs in the browser. Both talk to the API over HTTP(S), and both send the user to Keycloak to log in.

\section{Local Development}

\begin{center}
\begin{tabular}{@{}l l l p{5.5cm}@{}}
\textbf{Service} & \textbf{Image} & \textbf{Host port} & \textbf{Purpose} \\
\hline
\texttt{api} & built from \texttt{src/PrepMeUp.Api} & 5001 & ASP.NET Core Web API \\
\texttt{db} & SQL Server 2022 & none & Application data \\
\texttt{keycloak} & Keycloak 26.3 & 8082 & Login and roles, runs in dev mode \\
\end{tabular}
\end{center}

\begin{itemize}
    \item The database has no host port. Only the API can reach it, on \texttt{db:1433}.
    \item The API waits until the database health check passes before it starts.
    \item Keycloak imports the realm from \texttt{infra/keycloak/realm-prepmeup.json} at startup.
    \item The admin panel (Vite, port 5173) and the Expo app (port 8081) run outside Docker. These two origins are allowed by the API's CORS policy.
    \item Data is kept in the volumes \texttt{prepmeup-data} and \texttt{prepmeup-keycloak}.
\end{itemize}

\subsection{Token Validation}
The browser and the phone reach Keycloak on \texttt{localhost:8082}, but inside Docker the API reaches it on \texttt{keycloak:8080}. The token says it was issued by \texttt{localhost:8082}, so the API is configured with two addresses:
\begin{itemize}
    \item \texttt{Keycloak\_\_Authority}: the public address, which must match the issuer in the token.
    \item \texttt{Keycloak\_\_MetadataAddress}: the internal address, where the API fetches Keycloak's configuration and signing keys.
\end{itemize}

\section{Production}

\begin{center}
\begin{tabular}{@{}l l p{7cm}@{}}
\textbf{Service} & \textbf{Image} & \textbf{Purpose} \\
\hline
\texttt{api} & built from \texttt{src/PrepMeUp.Api} & ASP.NET Core Web API \\
\texttt{db} & SQL Server 2022 & Application data \\
\texttt{keycloak} & Keycloak 26.3 & Login and roles, runs in production mode \\
\texttt{keycloak-db} & PostgreSQL 17 & Keycloak's own data (users, sessions) \\
\end{tabular}
\end{center}

\begin{itemize}
    \item No container publishes a port. Coolify's reverse proxy receives the traffic from the internet on HTTPS and forwards it to \texttt{api} and \texttt{keycloak} on port 8080 inside the network.
    \item Because HTTPS ends at the proxy, Keycloak reads the \texttt{X-Forwarded} headers to know the original address and scheme.
    \item \texttt{db} and \texttt{keycloak-db} can only be reached from inside the network.
    \item The public URLs of the admin panel, the customer app and Keycloak are set as environment variables in Coolify. The same variables are used for the API's CORS policy and the redirect URIs in the Keycloak realm.
    \item All services have \texttt{restart: unless-stopped}, so they start again after a crash or a server restart.
\end{itemize}

\section{Login Flow}
Both clients use the Authorization Code flow with PKCE:
\begin{enumerate}
    \item The client opens Keycloak's login page in the browser.
    \item The user logs in at Keycloak. The password is never sent to the client or the API.
    \item Keycloak redirects back to the client with a short-lived code. For the app, the redirect uses the app scheme \texttt{prepmeup://}.
    \item The client exchanges the code for an access token (a JWT, valid for 15 minutes).
    \item The client sends the token to the API in the \texttt{Authorization: Bearer} header.
    \item The API checks the token's signature, issuer and audience (\texttt{prepmeup-api}), and uses the role (\texttt{household} or \texttt{admin}) for authorization.
\end{enumerate}

Keycloak already adds the audience \texttt{prepmeup-api} to the tokens for both clients. Step 6 is not implemented in the API yet.

\end{document}
